\subsection{Separation of convex sets in a topological vector space}

DEFINITION 1. --- Two non-empty sets A, B of a real topological vector space E are said to be separated by a closed hyperplane H if A is contained in one of the closed half-spaces determined by H and B is contained in the other closed half-space.

DEFINITION 2. --- Two non-empty sets A, B of a real topological vector space are said to be strictly separated by the closed hyperplane H if A is contained in one of the open half-spaces determined by H, and B is contained in the other open half-space.

PROPOSITION 1. --- Let A be an open non-empty convex set and let B be a non-empty convex set in a real topological vector space E; if A does not meet B then there exists a closed hyperplane that separates A from B.

For the set \(\mathbf{C} = \mathbf{A} - \mathbf{B}\) is open, convex (II, p.~9, prop. 7) and non-empty, also \(0 \notin \mathbf{C}\) . By theorem 1 of II, p.~36, there exists a continuous linear form \(f \neq 0\) on \(\mathbf{E}\) such that \(f(z) > 0\) in \(\mathbf{C}\) . Then, for all \(x \in \mathbf{A}\) , and \(y \in \mathbf{B}\) , we have \(f(x) > f(y)\) . Write \(\alpha = \inf_{x \in \mathbf{A}} f(x)\) ; \(\alpha\) is finite and we have \(f(x) \geqslant \alpha\) for all \(x \in \mathbf{A}\) and \(f(y) \leqslant \alpha\) for all \(y \in \mathbf{B}\) ; the closed hyperplane \(\mathbf{H}\) with the equation \(f(z) = \alpha\) separates \(\mathbf{A}\) from \(\mathbf{B}\) .

Remarks. --- 1) The hyperplane H does not meet A (II, p.~15, prop. 1; if A and B are two convex non-empty open sets that do not meet then there exists a closed hyperplane that separates A strictly from B.

2) However, when B is not open, it is not necessarily the case that there exists a closed hyperplane that separates A strictly from B, even if E is of finite dimension, and even if \(\overline{A}\) does not meet \(\overline{B}\) (II, p.~78, exerc. 12).

DEFINITION 3. --- For a subset A of a vector space E, a hyperplane H is called a support hyperplane of A, if H contains at least one point of A and all the points of A lie on the same side of H.

Let f be a linear form on E that is not identically zero; to say that the hyperplane of the equation \(f(x) = \alpha\) is a support hyperplane of A means that \(\alpha\) is either the smallest or the largest member of the set \(f(A) \subset R\) . In other words, there exists a support hyperplane of A parallel to the hyperplane of equation \(f(x) = 0\) , if, and only if, one of the bounds of the set \(f(A)\) is finite and belongs to \(f(A)\) .

PROPOSITION 2. --- Let A be a non-empty compact subset of a topological vector space E. For every closed hyperplane H in E, there exists a support hyperplane of A parallel to H.

For, if \(f(x) = \gamma\) is an equation of H, where f is a continuous linear form in E, the restriction of f to A is continuous, therefore bounded and attains its bounds in A (GT, IV, § 6.1, th. 1).

This demonstrates that there exist one or two support hyperplanes of A parallel to H; the first case can only arise when A is completely contained in a hyperplane parallel to H.

PROPOSITION 3. --- In a topological vector space E, let A be a closed convex set with a non-empty interior. Then every support hyperplane of \(\mathbf{A}\) is closed and every frontier point of \(\mathbf{A}\) belongs to at least one support hyperplane of \(\mathbf{A}\) .

Every support hyperplane of A is closed, since all the points of A are on the same side of the hyperplane (II, p.~15, prop. 17). Also if \(x_0\) is a frontier point of A, then \(x_0\) does not belong to the open non-empty convex set \(\mathring{\mathrm{A}}\) ; after th. 1 of II, p.~36 there exists a hyperplane H that contains \(x_0\) and does not meet \(\mathring{\mathrm{A}}\) . As A is the closure of \(\mathring{\mathrm{A}}\) (II, p.~14, cor. 1 to prop. 16), it follows from prop. 17 of II, p.~15 that H is a support hyperplane of A.

